\section{Moduli del sistema}
\label{sec:moduli}

Questa sezione descrive in dettaglio ciascun modulo software del progetto,
organizzato secondo i livelli logici già introdotti nella
Sezione~\ref{sec:architettura}: Sensor Layer, Edge AI Layer, Agent Layer,
Actuator Layer, e i due orchestratori (CLI e dashboard).

\subsection{Sensor Layer}

\subsubsection{Fotocamera virtuale (\texttt{sensors/virtual\_camera.py})}

La classe \texttt{VirtualCameraSensor} simula una fotocamera installata
nella serra, sostituendo un dispositivo fisico (collegabile via CSI o USB a
un Raspberry Pi reale) con un campionamento pseudo-casuale di immagini reali
dal dataset PlantVillage-Tomato. All'inizializzazione, il modulo scansiona
la cartella del dataset e costruisce un indice di tutte le immagini
disponibili, organizzato per classe. Ad ogni chiamata del metodo
\texttt{capture()}, viene selezionata un'immagine (in modo uniforme fra
tutte le classi, oppure con un \emph{bias} verso una classe specifica se
configurato tramite il parametro \texttt{bias\_class}, utile per simulare
scenari mirati come un'epidemia di una particolare patologia durante una
dimostrazione), viene applicato il preprocessing descritto nella
Sezione~\ref{sec:modello-ml}, e viene restituita una struttura dati
(\texttt{CameraCapture}) contenente il tensore pronto per l'inferenza, il
percorso del file, l'etichetta reale (nota perché proveniente da un dataset
etichettato, usata esclusivamente per il calcolo dell'accuratezza nel PoC e
mai fornita in input alla rete) e un timestamp.

\begin{notebox}
Per evitare che la CNN venga valutata su immagini già viste durante il
training (situazione che produrrebbe un'accuratezza artificiosamente vicina
al 100\%, non rappresentativa delle reali capacità del modello), la cartella
dataset utilizzata da PomodorIA contiene esclusivamente il sottoinsieme di
\emph{test} (20\% del dataset originale), ottenuto riproducendo esattamente
lo stesso split casuale (\texttt{random\_split} con seed fisso pari a 42)
utilizzato nel progetto di training originale. Questo garantisce che ogni
misura di accuratezza riportata in questo documento rifletta la reale
capacità di generalizzazione del modello su dati mai osservati in fase di
addestramento.
\end{notebox}

\subsubsection{Simulatore di sensori ambientali (\texttt{sensors/environment\_simulator.py})}

La classe \texttt{EnvironmentSensorSimulator} genera valori sintetici di
temperatura, umidità relativa, umidità del suolo e luminosità, in analogia
con sensori fisici reali (rispettivamente assimilabili a un DHT22, un
igrometro capacitivo e un sensore di luminosità BH1750). La particolarità
del modulo, e l'implementazione concreta del principio di \textbf{Data
Fusion} discusso nella Sezione~\ref{sec:contesto}, risiede nel metodo
\texttt{read(disease\_label)}: se viene fornita l'etichetta reale della
patologia rilevata nell'immagine appena acquisita, i valori ambientali
vengono generati secondo un profilo statistico specifico per quella
patologia (una distribuzione gaussiana con media e deviazione standard
calibrate su conoscenza botanica reale, ad esempio umidità elevata e
temperatura moderata per le patologie fungine come il Late Blight), anziché
in modo puramente casuale. Questo rende lo scenario simulato coerente e
"raccontabile": la diagnosi visiva e il contesto ambientale, per quanto
entrambi simulati, risultano tra loro plausibili, esattamente come
avverrebbe incrociando un dato reale da fotocamera con un dato reale da
sensore ambientale in un deployment fisico.

\subsection{Edge AI Layer}
\label{sec:edge-layer}

Il modulo \texttt{edge/inference\_engine.py} implementa la classe
\texttt{InferenceEngine}, responsabile di:

\begin{enumerate}
    \item \textbf{caricamento del modello}: in base alla configurazione
    (\texttt{full\_precision} oppure \texttt{optimized}), carica il
    checkpoint appropriato, ricostruendo l'architettura \texttt{TomatoCNN}
    e caricando i pesi salvati, oppure delegando l'esecuzione a un motore
    ONNX Runtime (Sezione~\ref{sec:compressione});
    \item \textbf{esecuzione dell'inferenza}: applica il modello al tensore
    fornito dalla fotocamera virtuale, misurando il tempo di esecuzione con
    precisione al millisecondo;
    \item \textbf{post-processing}: applica la funzione softmax ai logits
    grezzi restituiti dal modello, individua la classe con probabilità
    massima e restituisce una struttura dati (\texttt{InferenceResult})
    contenente etichetta predetta, confidenza, vettore di probabilità
    completo, tempo di inferenza e RAM utilizzata dal processo in quel
    momento;
    \item \textbf{simulazione dei vincoli di un dispositivo Edge}: il
    numero di thread CPU utilizzabili da PyTorch viene limitato a 1
    (\texttt{torch.set\_num\_threads(1)}) e l'esecuzione è forzata su CPU,
    per rendere le misure di latenza rappresentative di un dispositivo
    embedded privo di GPU dedicata, come un Raspberry Pi.
\end{enumerate}

Il motore di inferenza supporta inoltre la selezione esplicita, tramite il
parametro di configurazione \texttt{optimized\_variant}, di quale variante
compressa del modello utilizzare (dettagli in
Sezione~\ref{sec:compressione}), rendendo possibile confrontare
sperimentalmente le diverse tecniche di compressione applicate.

\subsection{Agent Layer}
\label{sec:agent-layer}

Il modulo \texttt{agent/decision\_agent.py} implementa l'agente decisionale,
il componente che realizza concretamente il modello \textbf{PEAS}
introdotto nella Sezione~\ref{sec:peas}.

\begin{table}[h]
\centering
\small
\begin{tabular}{@{}p{1.6cm}p{11.2cm}@{}}
\toprule
\textbf{PEAS} & \textbf{Definizione nel sistema} \\
\midrule
Performance & Diagnosi corrette, tempo di reazione, risparmio idrico, riduzione dei falsi allarmi. \\
Environment & Serra simulata: parzialmente osservabile (i sensori non colgono l'intero stato della pianta), dinamica, continua per i valori ambientali, discreta per le classi patologiche. \\
Actuators & Pompa di irrigazione, ventola di areazione, LED/allarme, notifica testuale. \\
Sensors & Fotocamera virtuale, sensori ambientali virtuali. \\
\bottomrule
\end{tabular}
\caption{Caratterizzazione PEAS dell'agente decisionale di PomodorIA.}
\label{tab:peas}
\end{table}

Dal punto di vista della tassonomia degli agenti, il \texttt{DecisionAgent}
si qualifica come \textbf{agente reattivo basato su modello e su obiettivi}:
mantiene uno stato interno (il contatore di rilevamenti di malattia
consecutivi, utile per distinguere un singolo rilevamento isolato da
un'epidemia in corso) e sceglie le proprie azioni tramite un insieme di
regole a priorità decrescente, riportate in Tabella~\ref{tab:regole-agente}.

\begin{table}[h]
\centering
\small
\begin{tabular}{@{}p{0.7cm}p{5.2cm}p{6.9cm}@{}}
\toprule
\textbf{\#} & \textbf{Condizione} & \textbf{Azione} \\
\midrule
1 & Confidenza della CNN inferiore alla soglia configurata & Allarme + notifica di richiesta ispezione umana \\
2 & Malattia virale rilevata & Allarme critico + notifica (possibile vettore da insetto) \\
3 & Malattia fungina rilevata + umidità sopra soglia & Attivazione ventilazione + notifica \\
4 & Malattia batterica o da parassiti rilevata & Notifica di livello informativo/warning \\
5 & Umidità del suolo sotto soglia (indipendente dalla diagnosi) & Attivazione irrigazione \\
6 & Temperatura sopra soglia (indipendente dalla diagnosi) & Attivazione ventilazione \\
7 & N rilevamenti di malattia consecutivi $\geq$ soglia & Allarme di "epidemia" simulata \\
8 & Pianta sana & Disattivazione di tutti gli attuatori di emergenza \\
\bottomrule
\end{tabular}
\caption{Regole condizione-azione dell'agente decisionale, in ordine di priorità.}
\label{tab:regole-agente}
\end{table}

Un aspetto rilevante da chiarire riguarda la regola 1: contrariamente a
un'interpretazione intuitiva, la richiesta di "ispezione umana" scatta
quando la confidenza della CNN è \emph{inferiore} alla soglia (di default
70\%, valore configurabile e scelto come convenzione di progetto, non
derivato da un'analisi statistica dei dati), non quando è superiore. La
soglia rappresenta quindi un meccanismo di cautela: se il modello non è
sufficientemente sicuro della propria predizione, il sistema segnala la
necessità di una verifica esterna, invece di agire automaticamente sulla
base di una diagnosi incerta.

\begin{warnbox}
Nell'implementazione attuale, la richiesta di "ispezione umana" è
\emph{simbolica}: si traduce nell'attivazione di un allarme visivo e nella
generazione di una notifica testuale, ma il ciclo automatico
sense-think-act-log prosegue comunque al passo successivo senza attendere
un effettivo intervento dell'operatore umano. Questo limite, e le possibili
estensioni per renderlo funzionalmente bloccante, sono discussi nella
Sezione~\ref{sec:limiti}.
\end{warnbox}

\subsection{Actuator Layer}

Il livello di attuazione è realizzato da due moduli distinti, che riflettono
la separazione fra hardware (simulato) e logica applicativa.

\subsubsection{Interfaccia GPIO simulata (\texttt{actuators/fake\_gpio.py})}

\textbf{GPIO} (\emph{General Purpose Input/Output}) indica i piedini fisici
di una scheda come il Raspberry Pi, utilizzabili via software per pilotare
dispositivi elettrici reali. Il modulo \texttt{fake\_gpio.py} replica
fedelmente l'interfaccia pubblica della libreria ufficiale \texttt{RPi.GPIO}
(le stesse costanti e le stesse funzioni: \texttt{setmode}, \texttt{setup},
\texttt{output}, \texttt{input}, \texttt{cleanup}), mantenendo tuttavia lo
stato dei "pin" esclusivamente in una struttura dati Python, senza pilotare
alcun hardware reale. Questa scelta implementativa rende il codice del
livello di attuazione immediatamente riutilizzabile su hardware reale:
sarebbe sufficiente sostituire l'istruzione di importazione del modulo
(\texttt{from actuators import fake\_gpio as GPIO}) con
\texttt{import RPi.GPIO as GPIO}, senza alcuna altra modifica.

\subsubsection{Banco attuatori (\texttt{actuators/actuators.py})}

Il modulo definisce una gerarchia di classi basata sull'ereditarietà, il
meccanismo standard della programmazione orientata agli oggetti per
condividere comportamento comune fra più classi specializzate senza
duplicare codice. La classe base \texttt{SimulatedActuator} implementa, una
sola volta, la logica di attivazione/disattivazione (gestione dello stato
interno, chiamata all'interfaccia GPIO, generazione di un log
strutturato dell'azione); le classi \texttt{IrrigationActuator},
\texttt{VentilationActuator} e \texttt{AlarmActuator} ereditano da essa,
specializzandola solo con un nome descrittivo e un numero di pin distinto,
senza reintrodurre alcuna logica propria. La classe
\texttt{NotificationActuator}, priva di un pin fisico associato, simula
l'invio di una notifica testuale con livello di gravità (informativo,
warning, critico). Infine, la classe \texttt{ActuatorBank} aggrega i quattro
attuatori in un unico punto di accesso, semplificando l'interazione da
parte dell'agente decisionale e fornendo un metodo di utilità per
disattivare tutti gli attuatori contemporaneamente (usato, ad esempio, alla
chiusura ordinata del programma).

\subsection{Orchestrazione: CLI e Dashboard}

\subsubsection{Orchestratore da riga di comando (\texttt{orchestrator/main\_loop.py})}

Il modulo implementa la classe \texttt{Orchestrator}, che istanzia tutti i
componenti di dominio descritti in questa sezione e ne coordina l'esecuzione
in un ciclo continuo, oltre alla classe \texttt{PoCLogger}, responsabile
della scrittura strutturata su file CSV di ogni ciclo eseguito (timestamp,
etichetta reale, predizione, confidenza, tempo di inferenza, valori
ambientali, stato degli attuatori). L'orchestratore gestisce inoltre in modo
esplicito l'interruzione da tastiera (CTRL+C), garantendo che tutti gli
attuatori vengano disattivati e il file di log chiuso correttamente prima
della terminazione del processo.

\subsubsection{Dashboard interattiva (\texttt{dashboard/app\_streamlit.py})}

La dashboard, realizzata con il framework \emph{Streamlit} e la libreria di
visualizzazione \emph{Plotly}, riproduce lo stesso identico ciclo
sense-think-act-log all'interno di un'interfaccia web, offrendo:

\begin{itemize}
    \item visualizzazione in tempo reale dell'immagine analizzata e delle
    probabilità softmax calcolate dalla CNN per ciascuna delle 10 classi;
    \item indicatori grafici (\emph{gauge chart}) per i valori dei sensori
    ambientali simulati, con andamento storico;
    \item indicatori di stato per ciascun attuatore, e visualizzazione del
    ragionamento dell'agente decisionale (le stringhe esplicative generate
    dalle regole condizione-azione);
    \item un pannello di monitoraggio delle risorse (CPU, RAM, numero di
    thread, latenza media di inferenza), a rappresentare i vincoli tipici
    di un dispositivo Edge;
    \item grafici di analisi aggregata: accuratezza cumulativa nel tempo,
    distribuzione delle predizioni, istogramma delle latenze, confronto fra
    le diverse varianti di compressione del modello.
\end{itemize}

Un controllo aggiuntivo, discusso nella Sezione~\ref{sec:compressione},
consente di selezionare manualmente quale specifica variante compressa del
modello utilizzare, in modo da poter dimostrare interattivamente
l'effettivo trade-off fra dimensione, velocità e accuratezza introdotto da
ciascuna tecnica.
